\subsection{Hard sectors and approximation loss}
\label{ian:hard-sector-definitions-forward}

Status: proved; all statements in this subsection hold for arbitrary bounded real-valued functions, without quantum realizability.

Let \(n\ge1\), let \(\Omega_n=\{-1,1\}^n\), and let
\(Q:\Omega_n\to[0,1]\).
For an integer \(t\ge0\), let \(\mathcal D_t\) denote the deterministic
coordinate decision trees of worst-case depth at most \(t\), with leaf
outputs in \([0,1]\), and put
\[
 \ell_\mu(t,Q)=\min_{D\in\mathcal D_t}\E_\mu|D-Q|.
\]
The minimum exists: there are finitely many tree structures and coordinate
labellings of bounded depth, and each associated space of leaf outputs
is compact.

Write \(\mathcal F_t\) for the corresponding trees with leaf outputs in
\([-1,1]\).

\begin{definition}[Hard laws and sectors]
\label{ian:hard-sector-definition}
For a probability law \(\nu\), put \(m_\nu=\E_\nu Q\).
For \(0<\gamma<a\le1/4\), the assertion
\(\operatorname{Hard}(Q,\nu;t,a,\gamma)\) means
\[
 \Var_\nu Q\ge a,
 \qquad
 \left|\E_\nu[(Q-m_\nu)F]\right|\le\gamma
 \quad\text{for every }F\in\mathcal F_t.
\]
A smooth sector of weight \(w\in(0,1]\) is such a law satisfying
\[
 w\nu(x)\le\mu(x)\qquad(x\in\Omega_n).
\]
A set sector of weight \(w\) is a set \(H\subseteq\Omega_n\) with
\(\mu(H)\ge w\) for which the conditional law \(\mu(\cdot\mid H)\)
is hard with the specified parameters.
Every set sector is a smooth sector of the same weight.
\end{definition}

\begin{lemma}[Aggregation and extremal tests]
\label{ian:hard-sector-tree-closure}
The following assertions hold for every \(n\ge1\) and integer \(t\ge0\).

First, any function of the outputs of \(k\) depth-\(t\) trees is
computable by a tree of depth at most \(kt\), provided the resulting
leaf outputs belong to the desired output range.

Second, for any law \(\nu\), checking
\[
 \left|\E_\nu[(Q-m_\nu)F]\right|\le\gamma
 \qquad(F\in\mathcal F_t)
\]
is equivalent to checking only trees with leaf outputs in
\(\{-1,1\}\).
The number of these Boolean-valued tests is at most
\[
 N(n,t)=n^{2^t-1}2^{2^t}.
\]
\end{lemma}

\begin{proof}
Run the \(k\) trees consecutively and then apply the prescribed function
to their outputs. Each run uses at most \(t\) queries. Previously
queried coordinates can be remembered.

For the second assertion, fix a query structure and let \(A_v\) be its
leaf events. If its leaf outputs are \(f_v\in[-1,1]\), then
\[
 \E_\nu[(Q-m_\nu)F]
 =\sum_v f_v\,\E_\nu[(Q-m_\nu)\mathbf1_{A_v}].
\]
The maximum absolute value over the cube of leaf assignments is attained
at an assignment with every \(f_v\in\{-1,1\}\).
Thus the Boolean-valued tests suffice.

A tree of depth at most \(t\) can be padded to the complete binary tree
of depth \(t\), using repeated queries if necessary. Its \(2^t-1\)
internal vertices each have at most \(n\) coordinate labels, and its
\(2^t\) leaves each have two output labels. This proves the count.
\end{proof}

\begin{proposition}[H1-forward]
\label{ian:hard-sector-forward}
For every \(n,\mu,Q,t,w,a,\gamma,\nu\) satisfying
Definition~\ref{ian:hard-sector-definition},
\[
 \nu\le\mu/w,\qquad
 \operatorname{Hard}(Q,\nu;t,a,\gamma)
 \quad\Longrightarrow\quad
 \ell_\mu(t,Q)\ge w(a-\gamma).
\]
\end{proposition}

\begin{proof}
Fix \(D\in\mathcal D_t\), and write \(m=\E_\nu Q\).
The tree \(F=2D-1\) belongs to \(\mathcal F_t\), so
\[
 2\Cov_\nu(Q,D)=\E_\nu[(Q-m)(2D-1)]\le\gamma.
\]
Consequently,
\[
 \E_\nu(Q-D)^2
 =\Var_\nu Q+\E_\nu(D-m)^2-2\Cov_\nu(Q,D)
 \ge a-\gamma.
\]
Because \(Q,D\in[0,1]\), one has \(|Q-D|\ge(Q-D)^2\).
Pointwise domination therefore gives
\[
 \E_\mu|Q-D|
 \ge w\E_\nu|Q-D|
 \ge w(a-\gamma).
\]
Taking the minimum over \(D\) proves the assertion. This is the
chosen-sense hard-law input needed by
Lemma~\ref{threshold:hard-component}.
\end{proof}
